\begin{claim}{D.2.51}{Proved}{The exact stable-spherical criterion and the three remaining orders}
Put $L=Y\langle7\rangle$ and retain the quadratic coordinates
\[
 \pi_{15}D_2L=\ZZ/8\{c\}\oplus\ZZ/2\{u\},
 \qquad h=h_Ml_4=\partial(2c),
 \tag{D.2.51a}
\]
with the compatible orientation of D.2.43.  For the actual
James--Hopf map of that claim,
\[
          \operatorname{im}\bigl(JH:\pi_{15}^sL
                      \longrightarrow\pi_{15}D_2L\bigr)
                         \subseteq\langle2c\rangle.
          \tag{D.2.51b}
\]
Let $a_7\in H^7(L;\ZZ/4)$, $b_8=\beta_4a_7$ and $\theta$
be the half-product class of D.2.45.  The following conditions
are equivalent:
\[
\begin{aligned}
 &h=0\text{ in }\pi_{14}L;\\
 &\text{there is }f\in\pi_{15}^sL\text{ with }
       \langle a_7\smile\operatorname{red}_4b_8,
                    \operatorname{red}_4h_{\ZZ}(f)\rangle=2;\\
 &\text{there is }f\in\pi_{15}^sL\text{ with }
       \langle\theta,\operatorname{red}_2h_{\ZZ}(f)\rangle=1.
\end{aligned}
\tag{D.2.51c}
\]
The last pairing is independent of the
$\operatorname{Sq}^1H^{14}(L;\FF_2)$ ambiguity of $\theta$.
More precisely, the possibilities are exactly
\[
\begin{array}{c|c|c}
 \operatorname{im}JH&\text{order of }h&
     \theta\text{ on the stable Hurewicz image}\\ \hline
 0&4&0\\
 \langle4c\rangle&2&0\\
 \langle2c\rangle&1&\text{nonzero}.
\end{array}
\tag{D.2.51d}
\]
Here order one means that $h$ is zero.  The integral class
constructed in D.2.50 satisfies the indicated pairing, but
its stable sphericity remains unproved.  Thus this criterion
does not choose a row of \textup{(D.2.51d)} or settle D.2.4.
\end{claim}
\begin{proof}
\textbf{The mod-two diagonal removes the independent coordinate.}
We use the full product vanishing in D.2.45j,
\[
                 H^7(L;\FF_2)\,H^8(L;\FF_2)=0,
                 \tag{D.2.51e}
\]
not only the vanishing of the particular reduction of
$a_7\operatorname{red}_4b_8$.  Since $L$ is six-connected,
the only summands of $H_{15}(L\wedge L;\FF_2)$ have
bidegrees $(7,8)$ and $(8,7)$.  Cup--diagonal duality and
\textup{(D.2.51e)} therefore imply
\[
 \overline\Delta_*\operatorname{red}_2z=0
             \quad\text{for every }z\in H_{15}(L;\ZZ).
 \tag{D.2.51f}
\]
All groups in these degrees are finite-dimensional over
$\FF_2$, so testing against the indicated cup products
detects the entire reduced diagonal.

To apply this to the actual stable operation, recall the
complete transfer computation in D.2.45f--h.  Its low
integral label chains are $a_7,b_8,r_8,s_9$, with
$db=4a$ and $ds=2r$.  They give
\[
 \begin{aligned}
 H_{15}(L\wedge L;\ZZ)
   &=\ZZ/4\{t=ab+ba\}\oplus\ZZ/2\{v=ar,w=ra\},\\
 \operatorname{tr}(c)&=t,\qquad
 \operatorname{tr}(u)=v+w.
 \end{aligned}
 \tag{D.2.51g}
\]
Here and below the transfer is followed by integral
Hurewicz.  Write $\bar a,\bar b,\bar r$ for the corresponding
mod-two homology classes.  The classes $\bar b,\bar r$
are independent in $H_8(L;\FF_2)$: the former is the
Moore-four Tor class and the latter is the reduction of
the integral order-two class.  Consequently the four
ordered tensors
\[
 \bar a\otimes\bar b,\quad \bar b\otimes\bar a,
 \quad \bar a\otimes\bar r,\quad\bar r\otimes\bar a
 \tag{D.2.51h}
\]
are linearly independent by the field K\"unneth isomorphism.
For $k\in\ZZ/8$ and $\epsilon\in\ZZ/2$, reduction of
\textup{(D.2.51g)} is zero precisely when
\[
             k\equiv0\pmod2,\qquad\epsilon=0.
             \tag{D.2.51i}
\]

For any stable sphere $f$, the transfer identity D.2.45e
gives
\[
       \operatorname{tr}JH(f)=-\overline\Delta_*h_{\ZZ}(f).
       \tag{D.2.51j}
\]
The sphere's reduced stable diagonal is zero.  The cited
identity is the actual Segal--Snaith James--Hopf operation,
with its comparison already checked in D.2.45 against
\cite[Theorems A(iii) and B]{Klein2026Hopf}.
Equations \textup{(D.2.51f)--(D.2.51j)} show that every
value $JH(f)=kc+\epsilon u$ has even $k$ and zero
$\epsilon$.  This proves \textup{(D.2.51b)} without
assuming anything about the stable Hurewicz image itself.

\textbf{An existential test, with the pointwise ambiguity retained.}
The coefficient functional
$a_7\otimes\operatorname{red}_4b_8$ evaluates as one on
$t$ and as zero on $v,w$.  If $JH(f)=kc$ with $k$ even,
\textup{(D.2.51j)} and D.2.45k give
\[
 \begin{aligned}
 \langle a_7\smile\operatorname{red}_4b_8,
                   \operatorname{red}_4h_{\ZZ}(f)\rangle
        &=-k\pmod4,\\
 2\langle\theta,\operatorname{red}_2h_{\ZZ}(f)\rangle
        &=-k\pmod4.
 \end{aligned}
 \tag{D.2.51k}
\]
The ambiguity of $\theta$ pairs to zero on an integral
reduction, since its homological first Bockstein is zero.
Thus the two pairing conditions in \textup{(D.2.51c)}
are equivalent, and each says exactly that
$JH(f)$ is $2c$ or $6c$.

The map $JH$ induces a homomorphism on the displayed
homotopy groups.  If the value is $6c$, replacing $f$ by
$3f$ changes it to $18c=2c$.  Hence existence of either
pairing is equivalent to $2c\in\operatorname{im}JH$.
The natural exact sequence of D.2.43 identifies the latter
condition with $\partial(2c)=0$, proving
\textup{(D.2.51c)}.  For a specified $f$, the pairing
still does not distinguish $2c$ from $6c$; the transfer
kernel $\langle4c\rangle$ from D.2.45 has not disappeared.
The new sufficiency is existential and uses the odd
rescaling of the sphere class.

Finally the image is a subgroup of the cyclic group
$\langle2c\rangle\cong\ZZ/4$, so it is one of
$0$, $\langle4c\rangle$ and $\langle2c\rangle$.
Exactness computes the order of the image of $2c$ in
each quotient as four, two and one, respectively.
Equation \textup{(D.2.51k)} computes the last column of
\textup{(D.2.51d)}.  This proves the complete table.
The exact finite certificate
\texttt{code/D51\_stable\_transfer\_range.py} checks all sixteen
quadratic-group elements, the mod-two range restriction, both
pairing-two values and their odd rescaling, and the three quotient
orders. It selects no actual James--Hopf image.
\end{proof}
